\chapter[Supplement G, Part 3: Latent Lie-Algebra Links]{Supplement G, Part 3:\\Latent Lie-Algebra Links}
\label{comp:gaussian-latent-links}

\section{The question and the answer}

\SG{ch:gauss-latent-links} evaluates the latent-link construction of
\cref{comp:lie-links} for multivariate Gaussian beliefs, the defining representation of
\(\GL^+(K)\), one belief channel, and held model latents. Random generators make links random
matrices. This changes what must be averaged, what must be normalized, and what a stationary link
means. The chapter gives explicit transported divergences and comparison-record currents,
evaluates expected loop defects under product recognition, integrates undirected Gaussian matter
out exactly, and grounds the linearized history equation in a ring spectrum. It also separates
two questions that fixed-link calculations cannot answer: whether the joint generator--matter
law is normalizable, and whether an unregularized optimum belongs to the exponential chart.

The earlier Gaussian chapters remain conditional calculations at a given link family. They
describe \(\Lambda(\mathring\omega)\), its Schur complements, and its Gaussian laws after the
generators have been fixed. Integrating over generators generally produces a mixture, which
need not be Gaussian (\MT{prop:link-conditional-transfer}). This is not a single Gaussian
evaluated at a mean link.

\section{The picture}

\begin{analogy}{exact}
Undirected Gaussian matter is a quadratic field coupled to link matrices. Integrating it out
gives a fluctuation determinant and, with an information vector, a source term mediated by the
propagator. Its convergent hopping expansion sums weighted closed walks
(\SG{prop:glink-induced-action}, \SG{prop:glink-hopping}). This is the Gaussian form of an induced
gauge action \citep{KazakovMigdal1993}. The matter lives in belief fibers over the agent network
at one base point; no base-manifold transport has been introduced.
\end{analogy}

Two energies must be kept distinct. A transported relative entropy compares whole belief laws.
A comparison record compares latent samples through a conditional observation kernel. A
quadratic mean mismatch with $W_e=\Sigma_i^{-1}$ equals the transported divergence when the
transported source covariance equals the receiver covariance
(\SG{prop:glink-law-level-energy}). Averaging a sample mismatch adds variances;
discarding them changes the force on the link.

\begin{analogy}{structural}
The ring grounds the Proca-type operator of \cref{comp:lie-links}: the anchor supplies positive
local stiffness, a cycle defect penalizes one collective edge mode, matter supplies a declared
isotropic stiffness, and increments penalize differences between episodes. The episode term has
Euclidean sign. The homogeneous most-probable-history recursion has increasing and decreasing
real solutions, rather than oscillations (\SG{prop:glink-linear-spectrum}). This is a history
posterior equation, not a real-time wave equation.
\end{analogy}

\section{Objects and notation}

\begin{decoder}
\(\mathring\omega_e\), \(\mathbf M=e^{\mathring\omega_e}\) & Reference-coordinate generator and represented link; reversal negates the generator & Link coordinate and transporter\\
\(\mu_i,\Sigma_i\) & Matter mean and covariance in reference coordinates & Field mean and propagator\\
\(\mathring\mu_e,\mathring\Lambda_e^{-1}\) & Mean and covariance of a generator recognition law & Link-field mean and fluctuations\\
\(\mathbf L(\mathbf B,\mathbf T)\) & Fr\'echet derivative of the exponential in direction \(\mathbf T\) & Jacobian from generator to matrix\\
\(\mathbf J_e\) & Negative matrix derivative of expected record energy & Prediction-error current\\
\(\mathfrak Q_p\) & Half the squared metric defect of a loop product from identity & Curvature penalty\\
\(\mathscr S^{\mathrm{ind}}=-\mathsf A\) & Negative matter log partition at the links & Induced effective action\\
\(\Lambda^{\mathrm{diag}},\Lambda^{\mathrm{hop}}\) & Matrices satisfying \(\Lambda=\Lambda^{\mathrm{diag}}-\Lambda^{\mathrm{hop}}\) & Local inverse propagator and hopping operator\\
\end{decoder}

Re-expressing local frames puts the matter laws in reference coordinates and retains the metric
\(\mathsf G\) in contractions. This does not remove independent twists. A ``chart link'' belongs
to \(\exp\mathfrak{gl}(K,\R)\), a proper subset of \(\GL^+(K)\): positive determinant does not
guarantee a real logarithm (\MT{prop:link-chart-image}). A critical point of the exponential
means its derivative is singular, not that an energy is minimized.

\clearpage
\section{The formulas, explained}

\subsection{Law divergence and sample energy}

For nondegenerate Gaussian laws \(q_i=\mathcal N(\mu_i,\Sigma_i)\) and
\(q_j=\mathcal N(\mu_j,\Sigma_j)\),
\begin{keyequation}{The transported Gaussian divergence}
\[
\begin{aligned}
\KL(q_i\Vert e^{\mathring\omega}_\#q_j)
={}&\tfrac12\Tr\bigl(\Sigma_j^{-1}e^{-\mathring\omega}\Sigma_i e^{-\mathring\omega^\top}\bigr)\\
&+\tfrac12(e^{-\mathring\omega}\mu_i-\mu_j)^\top\Sigma_j^{-1}(e^{-\mathring\omega}\mu_i-\mu_j)\\
&-\tfrac K2+\tfrac12\log\frac{\det\Sigma_j}{\det\Sigma_i}+\Tr\mathring\omega.
\end{aligned}
\]
\end{keyequation}
This is \SGeq{eq:glink-law-level-energy}, proved in \SG{prop:glink-law-level-energy}. The first
term compares covariances after pulling the receiver back through the link, the second compares
means in the source precision, and the determinant terms compare volume.
\(\Tr\mathring\omega\) is the link's log-volume change. At zero generator this is ordinary
Gaussian relative entropy. When
\(\Sigma_i=e^{\mathring\omega}\Sigma_j e^{\mathring\omega^\top}\), the covariance and volume terms
cancel and only the quadratic mean mismatch remains. For equal laws the local second-order form
is the orbit Fisher form of \SGeq{eq:glink-orbit-fisher}; stabilizer directions can have zero
stiffness (\MT{prop:link-law-level-energy}).

The comparison kernel instead reads
\(o_e\given z_i,z_j\sim\mathcal N(z_i-\mathbf Mz_j,W_e^{-1})\). Product matter recognition gives
\begin{keyequation}{Expected record energy and link current}
\[
\begin{aligned}
\E\mathsf E_e
={}&\tfrac12(\mu_i-o_e-\mathbf M\mu_j)^\top W_e(\mu_i-o_e-\mathbf M\mu_j)
+\tfrac12\Tr(W_e\Sigma_i)\\
&+\tfrac12\Tr(W_e\mathbf M\Sigma_j\mathbf M^\top)
+\tfrac12\log\det(2\pi W_e^{-1}),\\
\mathbf J_e={}&W_e\bigl[(\mu_i-o_e-\mathbf M\mu_j)\mu_j^\top-\mathbf M\Sigma_j\bigr],
\qquad \nabla_{\mathbf M}\E\mathsf E_e=-\mathbf J_e.
\end{aligned}
\]
\end{keyequation}
These are \SGeq{eq:glink-expected-record} and \SGeq{eq:glink-current}. Receiver variance is a
link-independent offset. Transported source variance changes with the link and contributes to
its current. The generator gradient is \(-\mathbf L(\mathring\omega_e^\top,\mathbf J_e)\),
the current pulled back through the exponential Jacobian (\SG{prop:glink-sample-record}). If
generators are random and independent of matter, the energy uses both
\(\E e^{\mathring\omega_e}\) and
\(\E[e^{\mathring\omega_e}\otimes e^{\mathring\omega_e}]\); replacing them by the exponential of
the generator mean drops fluctuations.

Without a prior the unique minimizer over real transport matrices is
\[
\mathbf M^\star=(\mu_i-o_e)\mu_j^\top(\Sigma_j+\mu_j\mu_j^\top)^{-1}.
\]
Its rank is at most one, so for \(K\geq2\) it lies outside the invertible group. This uses
independent matter recognition: cross covariance replaces the numerator by a second-moment
regression, which can be invertible (\MT{cor:link-prior-necessary}). This is not a universal
obstruction to learned links.

For $K\geq2$ the exponential chart adds a distinction. If \(\Tr\mathbf M^\star\geq0\), its infimum equals
the unrestricted minimum and is not attained. If \(\Tr\mathbf M^\star<0\), the chart infimum is
strictly higher: no sequence of real exponentials approaches \(\mathbf M^\star\). In this second
case nonattainment is proved at every chart link whose negative eigenvalues are semisimple,
covering all chart links for \(K=2,3\) (\SG{prop:glink-chart-infimum-semisimple}). Attainment at
nonsemisimple negative eigenvalues for \(K\geq4\) is open. Stationary generators can nevertheless
exist where the exponential derivative is singular; the two-dimensional example of
\SG{prop:glink-sample-record}(iv) is a saddle. Stationarity is not an optimum.

\subsection{Fluctuating holonomy}

Take \(R=\bigotimes_{e\in E_+}\mathcal N(\mathring\mu_e,\mathring\Lambda_e^{-1})\) and a closed
walk \(p=(e_0,\ldots,e_{r-1})\) using each reversal pair at most once. Write
\(\mathbf H_a=\E_R e^{\mathring\omega_{e_a}}\) and
\(\mathbf T_a=\E_R[e^{\mathring\omega_{e_a}}\otimes e^{\mathring\omega_{e_a}}]\).
\begin{keyequation}{The expected metric loop defect}
\[
\E_R\mathfrak Q_p
=\tfrac12(\mathbf G^+)^\top\left(\prod_{a=0}^{r-1}\mathbf T_a\right)\mathbf G^-
-\Tr\left(\prod_{a=0}^{r-1}\mathbf H_a\right)+\tfrac K2.
\]
\end{keyequation}
Here \(\mathbf G^+_{(a,b)}=\mathsf G_{ab}\),
\(\mathbf G^-_{(a,b)}=(\mathsf G^{-1})_{ab}\), and both products preserve the walk's order
(\SG{prop:glink-expected-loop}, \SGeq{eq:glink-expected-defect}). Independence gives products of
moments, not exponentials of means. Correlated recognition needs joint moments. Even product
recognition does not permit this factorization when a reversal pair repeats: a scalar backtrack
has holonomy identically one, while multiplying its separate exponential means gives
\(e^{\mathring\Lambda_e^{-1}}>1\).

At zero means and small covariances,
\(\E_R\mathfrak Q_p=\tfrac12\sum_a\Tr\mathring\Lambda_{e_a}^{-1}
+O(\max_a\lVert\mathring\Lambda_{e_a}^{-1}\rVert^2)\).
Flat means therefore do not make a fluctuating loop flat. At nonzero means the correction also
uses a second derivative of the exponential (\SGeq{eq:glink-defect-first-order}).

\subsection{Integrating out matter and checking the joint law}

Undirected matter has the energy
\[
\begin{aligned}
\mathsf E(z;\mathring\omega)
={}&\tfrac12\sum_i z_i^\top A_i z_i-\mathsf h^\top z\\
&+\tfrac12\sum_{e:j\to i\in E_+}(z_i-e^{\mathring\omega_e}z_j)^\top W_e(z_i-e^{\mathring\omega_e}z_j).
\end{aligned}
\]
Its precision is \(\Lambda(\mathring\omega)=A_\oplus+L(\mathring\omega)\). With every
\(A_i\succ0\), integration gives
\begin{keyequation}{The action induced by Gaussian matter}
\[
\begin{aligned}
\mathscr S^{\mathrm{ind}}(\mathring\omega)
={}&\tfrac12\log\det\Lambda(\mathring\omega)
-\tfrac12\mathsf h^\top\Lambda(\mathring\omega)^{-1}\mathsf h-\tfrac{NK}{2}\log2\pi,\\
P^{\mathscr T}(d\mathring\omega)
={}&\mathfrak Z^{-1}e^{-\mathscr S(\mathring\omega)-\mathscr S^{\mathrm{ind}}(\mathring\omega)}
\nu^{\mathscr T}(d\mathring\omega).
\end{aligned}
\]
\end{keyequation}
The action is minus the matter log partition, so the positive half-log-determinant accompanies
a negative source energy (\SGeq{eq:glink-induced-action}). The bound
\(\Lambda\succeq A_\oplus\succ0\) makes the whole joint normalizer finite for every finite base
measure after a normalizable prior tilt (\SG{prop:glink-induced-action}). The induced action alone
is not integrable against generator Lebesgue measure when an edge is present. Matter integration
does not replace the generator prior.

This whole-joint normalization is a distinct model choice. In a directed, conditionally normalized
record construction, integrating matter and records gives unit mass at each link family; it
induces no prior action of this kind, although observed records update the link posterior
(\MT{prop:link-induced-action}). In the undirected model the earlier Schur scalar and pair
normalizers are values of the displayed action. A link-dependent normalizer cannot be dropped
during link inference.

For noncompact links the action need not be purely a sum of loop defects. On a tree it can vary
through transported self terms and weights. It is constant there with isotropic self terms and
weights, zero information vector, and isometric links. The one-edge calculation of
\SGeq{eq:glink-edge-mass} gives stiffness along symmetric directions and a linear trace-direction
term even without a loop.

On an unanchored cycle,
\[
\det\Lambda(\mathring\omega)=\det(I-\mathring\Theta_p)^2\prod_{e\in p}\det W_e.
\]
The matter law exists at fixed links away from a null set, but its partition function diverges
at least as \(\lvert\det(I-\mathring\Theta_p)\rvert^{-1}\) across a regular singular hypersurface. A
finite base measure with positive density near such a point, including any nondegenerate Gaussian
base law, gives an infinite joint normalizer (\SG{cor:glink-unanchored-cycle}). This result is
about a cycle, not every unanchored graph. At \(K=1\), a connected multigraph with positive
weights, zero self terms, zero information vector, and nondegenerate Gaussian generator base law
is jointly normalizable when its cycle rank is at least two, and is not when its cycle rank is
one (\SG{prop:glink-unanchored-scalar}). The matrix question for \(K\geq2\) remains open.

\subsection{Hopping and the Wilson triangle}

At zero information vector with positive-definite self terms, set
\(\mathbf H=(\Lambda^{\mathrm{diag}})^{-1}\Lambda^{\mathrm{hop}}\). Its spectral radius is below
one, and the absolutely convergent expansion is
\[
\mathscr S^{\mathrm{ind}}=\tfrac12\log\det\Lambda^{\mathrm{diag}}
-\tfrac12\sum_{r\geq1}\frac{\Tr\mathbf H^r}{r}-\tfrac{NK}{2}\log2\pi.
\]
Each trace sums length-\(r\) closed walks with a weight at every vertex
(\SG{prop:glink-hopping}). In general these are weighted matrix products. Isotropic data and
isometric links give scalar hopping ratios and a link-independent diagonal determinant. On the
isolated triangle in that setting,
\begin{keyequation}{The leading induced loop term on a triangle}
\[
\mathscr S^{\mathrm{ind}}(\mathring\omega)-\mathscr S^{\mathrm{ind}}(0)
=\frac{w_{e_0}w_{e_1}w_{e_2}}
{a^{\mathrm{diag}}_{i_0}a^{\mathrm{diag}}_{i_1}a^{\mathrm{diag}}_{i_2}}
\mathfrak Q_p(\mathring\omega)+\mathfrak o(\mathring\omega).
\]
\end{keyequation}
Here \(\mathfrak Q_p=K-\Tr\mathring\Theta_p\) and the remainder obeys
\[
\lvert\mathfrak o\rvert\leq\frac K5\left(
\frac{(2w^{\mathrm{hop}})^5}{1-2w^{\mathrm{hop}}}
+\frac{2(w^{\mathrm{hop}})^5}{1-w^{\mathrm{hop}}}\right),
\qquad w^{\mathrm{hop}}=\frac{\max_e w_e}{\min_i a_i^{\mathrm{diag}}}<\tfrac12.
\]
This is \SGeq{eq:glink-wilson-triangle}. The small parameter controls hopping, not an assumption
that each generator is small. The first link-dependent term is a positive loop penalty. Equality
of the full induced action with the chosen prior loop action is not proved. The general walk
expansion needs no isometry; this triangle reduction and bound do.

\subsection{The explicit ring spectrum}

For \(N\geq2\) ring agents and \(n_1+1\) episodes, let anchor precision \(\mathrm a>0\),
declared matter stiffness \(\mathrm b\geq0\), loop coefficient \(\mathrm c\geq0\), and increment
precision \(\mathbf m\geq0\) be isotropic. The cycle map sums edge generators. The Hessian is
\[
\mathbf K=\left[(\mathrm a+\mathrm b)I_N\otimes I_{n_1+1}
+\mathrm c\mathbf1\mathbf1^\top\otimes I_{n_1+1}
+\mathbf m I_N\otimes L_{n_1}\right]\otimes\mathrm{Id}_{\mathfrak g}.
\]
The cycle term has rank one in the edge index, rather than being the ring graph Laplacian. It
penalizes total generator around the ring and leaves zero-edge-sum profiles unchanged. The episode
term is the path Laplacian including endpoint rows. The eigenvalues are
\[
\begin{gathered}
\mathrm a+\mathrm b+\mathrm cN+4\mathbf m\sin^2\frac{\pi l}{2(n_1+1)},\qquad
\mathrm a+\mathrm b+4\mathbf m\sin^2\frac{\pi l}{2(n_1+1)},\\
l=0,\ldots,n_1.
\end{gathered}
\]
For each episode mode the first acts on the edge-uniform sector of dimension \(K^2\), the second
on the zero-edge-sum sector of dimension \((N-1)K^2\). Dimensions add when values coincide.
Every eigenvalue is at least \(\mathrm a+\mathrm b>0\) (\SG{prop:glink-linear-spectrum}). For
\(\mathbf m>0\) the homogeneous episode recursion has positive reciprocal roots. These signs
ground the Euclidean reading, but prove no nonlinear recognition-flow convergence.

\section{What is proved, and how solid it is}

\begin{scorecard}
\MT{def:link-exponential-chart} & Independent generators are latent link coordinates & \CHOICE & Exponential chart at one base point\\
\SG{prop:glink-law-level-energy} & Exact Gaussian divergence and generator derivative & \SOLID & Nondegenerate laws, \(\GL^+(K)\)\\
\SG{prop:glink-sample-record} & Expected energy, covariance current, rank-one optimum & \COND & Product matter recognition; critical chart stationarity may be a saddle\\
\SG{prop:glink-chart-infimum-semisimple} & No minimum at semisimple negative chart eigenvalues & \SOLID & Negative-trace case; covers all chart links for \(K=2,3\)\\
\SG{prop:glink-expected-loop} & Ordered exponential moments give expected defect & \COND & Product generator recognition, distinct reversal pairs\\
\SG{prop:glink-induced-action} & Exact action and finite joint normalizer & \SOLID & Every self term positive definite; induced action alone insufficient\\
\SG{cor:glink-unanchored-cycle} & Conditional law almost everywhere need not yield joint law & \SOLID & Unanchored cycle; density positive near a regular singularity\\
\SG{prop:glink-unanchored-scalar} & Cycle rank at least two restores normalization & \TOY & Scalar fibers, zero information vector, Gaussian base law\\
\SG{prop:glink-hopping} & Convergent weighted closed-walk expansion & \SOLID & Positive-definite self terms, zero information vector\\
\SG{prop:glink-hopping}\newline (v) & Leading triangle term is metric loop defect & \TOY & Isotropic isometric triangle, stated hopping bound\\
\SG{prop:glink-linear-spectrum} & Positive ring-history operator; nonoscillatory homogeneous recursion & \TOY & One cycle, isotropic coefficients\\
\SG{open:glink-unanchored-normalization} & Unanchored matrix-network joint normalization & \OPEN & \(K\geq2\), cycle rank at least two\\
\SG{open:glink-chart-infimum} & Attainment at nonsemisimple negative eigenvalues & \OPEN & \(K\geq4\), negative-trace regression optimum\\
\end{scorecard}

The conditional badges identify substantive recognition restrictions, not gaps in the Gaussian
calculations. Product matter recognition is what forces rank-one regression; correlations can
change the conclusion. Chart nonattainment is general for semisimple negative eigenvalues, while
its \(K=2,3\) consequences do not settle the remaining Jordan cases.

\begin{hygiene}
Conditions that matter: the real defining representation of \(\GL^+(K)\), nondegenerate matter
and generator covariances, product recognition where used, distinct reversal pairs for products
of loop moments, positive-definite self terms for general anchored normalization and hopping,
and joint normalization of undirected matter. The scalar unanchored theorem needs zero
information vector; the triangle needs isotropic data, isometric links, and the stated hopping
bound; the ring needs its isotropic coefficients. These are not results for arbitrary noncompact
nonabelian networks. Hygiene: finite graph inventories, Gaussian moments, reference measures, and
the matrix derivative conventions of \SG{sec:glink-setting}.
\end{hygiene}

\section{Toy versus general}

The energies, current, anchored induced action, and hopping expansion apply at arbitrary finite
\(N,K\) within their hypotheses, including 512 or more agents. The scalar cycle-rank theorem
does not settle the matrix problem. The triangle identifies how a positive loop penalty can arise
from integrating matter out under its isometry and hopping conditions; it is not proved to
dominate arbitrary large noncompact networks. The ring supplies a spectrum at any finite \(N\),
but its single cycle term and isotropic matter stiffness are special. Its positive finite gap
does not exclude collective behavior, critical slowing down, or symmetry breaking in other
large-network or limiting regimes.

\section{Open ends}

\SG{open:glink-unanchored-normalization} asks whether
\(\int\nu(d\mathring\omega)\det\Lambda(\mathring\omega)^{-1/2}\) is finite for connected matrix
multigraphs of cycle rank at least two, positive weights, zero self terms, zero information vector,
and product nondegenerate Gaussian generator base laws. A proof for every such graph and law at a
stated \(K\geq2\), or a divergent example in that class, settles the question. Almost-sure
positive definiteness of conditional precision alone does not settle it.

\SG{open:glink-chart-infimum} concerns negative-trace regression at \(K\geq4\). The unresolved
candidates have a negative eigenvalue with Jordan blocks of size at least two, paired as required
for a real logarithm. Excluding every such candidate as a chart minimizer, or exhibiting one that
attains the infimum for admitted data, settles it. The nonnegative-trace case and the \(K=2,3\)
cases are already closed negative results. Neither open item blocks the anchored latent-link
model. Its coarse closure is the separate question \MT{open:link-coarse-closure}, discussed in
\cref{comp:lie-links}.

\begin{verdictbox}{\DELIVERED}
\textbf{Claimed.} Gaussian values of transported law divergence, record current and unregularized
chart behavior, expected loop defects under product recognition, the action induced by undirected
matter and its hopping expansion, and a ring spectrum, with matrix unanchored normalization and
the remaining chart-attainment cases left open.

\textbf{Proved.} These calculations and negative results are provided at their stated generality.
Gaussian integration and the log-determinant expansion are standard in substance; their role here
is to retain the link-dependent normalizer and distinguish a conditional Gaussian law from the
whole joint model.

\textbf{Gaps.} No nonlinear convergence, anchor or covariance selection, equality of the full
induced action with the prior loop action, or large-network limit is proved. The scalar, triangle,
and ring results ground mechanisms in their named regimes. The two matrix questions and
latent-link coarse closure remain open.
\end{verdictbox}
